\section{Исправление функций вознаграждения на основе человеческих предпочтений для предотвращения злоупотребления вознаграждением}

\subsection{Ключевые идеи}

Статья \cite{hatgiskessell2025pbr} предлагает метод 
\emph{метод исправления наград на основе предпочтений 
(Preference-Based Reward Repair, PBRR)}, цель которого ~-- исправлять ошибочную 
прокси-функцию вознаграждения посредством предпочтений человека, а не обучать 
новую награду с нуля.

Основные идеи:

\begin{itemize}
    \item человек способен задать приблизительную прокси-награду 
    $\hat r_{\text{proxy}}$, которая содержит лишь локальные ошибки;
    \item можно корректировать только те переходы, где прокси-награда приводит к 
          неправильному поведению агента;
    \item предпочтения сравниваются между траекториями, полученными от 
          политики, оптимизирующей исправленную награду, и ссылочной политики;
    \item вводится специальная функция потерь, корректирующая только ошибочно 
    высокие значения награды.
\end{itemize}

\subsection{Математическое описание метода}

Метод \emph{исправления наград на основе предпочтений 
(Preference-Based Reward Repair, PBRR)} опирается на идею того, что вручную 
заданная человеком прокси-функция вознаграждения $\hat r_{\text{proxy}}$ 
в большинстве случаев верно отражает намерения человека, но содержит 
локальные ошибки, которые приводят к искажению вознаграждения. 
Вместо того чтобы обучать новую функцию вознаграждения с нуля, 
метод PBRR вносит в прокси-награду лишь минимальные целевые корректировки, 
используя специально сконструированный \emph{корректирующий член функции 
вознаграждения} $g_t(s,a,s')$, который обучается по человеческим предпочтениям.

Итоговая модифицированная функция награды на итерации $t$ задаётся аддитивно:
\begin{equation}
    \hat r_t(s,a,s') = \hat r_{\text{proxy}}(s,a,s') + g_t(s,a,s'),
\end{equation}
где $g_t$ параметризуется нейросетевой моделью и изменяет только те переходы, 
на которых прокси-награда приводит к нежелательному поведению агента. 
Таким образом, большая часть структуры $\hat r_{\text{proxy}}$ сохраняется 
неизменной, а исправления касаются лишь локальных ошибочных участков 
пространства состояний и действий.

Соответствие предпочтений человека и ранжирования траекторий, формируемого 
текущей наградой $\hat r_t$, моделируется \emph{моделью предпочтений Bradley--Terry}, 
которая задаёт вероятность того, что человек предпочитает траекторию 
$\tau_1$ траектории $\tau_2$, в виде
\begin{equation}
    P(\tau_1 \succ \tau_2 \mid \hat r) = 
    \sigma\!\left( 
        \hat r(\tau_1) - \hat r(\tau_2)
    \right),
    \qquad 
    \sigma(x) = \frac{1}{1 + e^{-x}}.
\end{equation}
Здесь $\hat r(\tau)$ представляет собой суммарное вознаграждение траектории, 
а сигмоида $\sigma(\cdot)$ задаёт гладкую зависимость вероятности от разницы 
вознаграждений.

Ключевое нововведение метода заключается во введении специальной функции потерь, 
которая принуждает корректирующий член $g$ изменять только действительно 
ошибочные части прокси-награды. Для этого все пары предпочтений разделяются 
на два множества: $D_t^+$ и $D_t^-$. В множество $D_t^+$ входят пары траекторий, 
для которых прокси-награда ранжирует их в соответствии с предпочтением человека, 
поэтому значения $g$ на таких траекториях должны оставаться близкими к нулю. 
В множество $D_t^-$ включаются пары, где прокси-награда ошибается, 
и именно эти переходы требуют корректировки. 

Функция потерь имеет вид
\begin{equation}
\begin{aligned}
L(g; \hat r_{\text{proxy}}, D_t) 
 =\,& L_{\text{pref}}(\hat r_{\text{proxy}} + g; D_t) \\
 &+ \lambda_1 \frac{1}{|D_t^+|}
     \sum_{(\tau_1, \tau_2) \in D_t^+}
     \left(g(\tau_1)^2 + g(\tau_2)^2\right) \\
 &+ \lambda_2 \frac{1}{|D_t^-|}
     \sum_{(\tau_1, \tau_2) \in D_t^-}
        \left[
            \mathbf{1}\{\tau_1 \succ \tau_2\} \, g(\tau_1)^2
            +
            \mathbf{1}\{\tau_2 \succ \tau_1\} \, g(\tau_2)^2
        \right].
\end{aligned}
\end{equation}

Первый член $L_{\text{pref}}$ представляет собой стандартную кросс-энтропию, 
используемую для согласования награды с человеческими предпочтениями. 
Однако он не различает корректные и ошибочные участки прокси-награды. 

Второй член, пропорциональный $\lambda_1$, подавляет корректирующий сигнал 
$g$ на тех траекториях, которые уже правильно оценены прокси-наградой. 
Это предотвращает разрушение исходной структуры $\hat r_{\text{proxy}}$ и 
обеспечивает стабильность метода. 

Третий член, пропорциональный $\lambda_2$, изменяет только те переходы, 
на которых прокси-награда противоречит предпочтениям человека, причём 
корректировка направлена на уменьшение награды ошибочно предпочтительных 
траекторий. Такая асимметрия исправлений делает метод устойчивым: вместо того 
чтобы повышать вознаграждение хороших траекторий, что может нарушить уже 
согласованные стимулы, метод понижает вознаграждение неверно оценённых переходов, 
устраняя стимулы, приводящие к \emph{искажению вознаграждения}.

В совокупности эти механизмы обеспечивают точечное исправление только тех 
элементов функции вознаграждения, которые действительно ответственны за 
некорректное поведение агента, при этом сохраняя структуру и полезные 
свойства исходной прокси-награды. Это позволяет устранять 
локальные и системные ошибки в ручной спецификации вознаграждения и 
предотвращать появление новых форм искажения вознаграждения.


\subsection{Алгоритм}

Алгоритм исправления наград на основе предпочтений реализуется как итеративная 
схема обновления корректирующего сигнала $g_t$, добавляемого к исходной 
прокси-функции вознаграждения $\hat r_{\text{proxy}}$. На каждой 
итерации агент обучается по текущей модифицированной награде $\hat r_t$, 
после чего его поведение сравнивается с поведением 
\emph{ссылочной политики} $\pi_{\text{ref}}$. Выявленные различия 
преобразуются в предпочтения человека между парами траекторий и используются 
для обновления корректирующего члена.

Алгоритм начинается с инициализации $g_0 = 0$, что задаёт начальную 
модифицированную награду $\hat r_0 = \hat r_{\text{proxy}}$. Далее 
итеративно выполняются следующие шаги:

\begin{enumerate}
    \item Устанавливается $g_0 = 0$ и вычисляется начальная функция 
    награды $\hat r_0 = \hat r_{\text{proxy}}$.

    \item На итерации $t$ обучается политика $\pi_t$, оптимизирующая 
    текущую награду $\hat r_t$. Для обучения используется метод 
    \emph{оптимизации проксимальных политик}, обеспечивающий стабильное 
    обновление параметров.

    \item Генерируются траектории из $\pi_t$ и ссылочной политики 
    $\pi_{\text{ref}}$. Эти траектории используются для выявления расхождений 
    между поведением, обусловленным прокси-награды, и желаемым 
    человеческим поведением.

    \item Полученные траектории попарно сравниваются человеком. Для каждой 
    пары $(\tau_1, \tau_2)$ фиксируется предпочтение, после чего пары 
    распределяются в множества $D_t^+$ и $D_t^-$ в зависимости от того, 
    согласуется ли ранжирование траекторий, заданное прокси-наградой, 
    с человеческим предпочтением.

    \item На основе множеств $D_t^+$ и $D_t^-$ обновляется корректирующий 
    член $g_{t+1}$ путём минимизации функции потерь
    $
        L(g; \hat r_{\text{proxy}}, D_t),
    $
    описанной в разделе ``Математическое описание метода''. 
    Обновление затрагивает лишь те переходы, которые ответственны за 
    несогласие с предпочтениями человека.

    \item Модифицированная функция вознаграждения обновляется по правилу
    \begin{equation}
        \hat r_{t+1}(s,a,s') = \hat r_{\text{proxy}}(s,a,s') + g_{t+1}(s,a,s'),
    \end{equation}
    после чего процедура повторяется на следующей итерации.
\end{enumerate}

Итерационная структура алгоритма обеспечивает локальное исправление ошибок 
в прокси-награде при сохранении её корректных свойств.


\subsection{Наборы данных}

В экспериментах используются четыре среды, каждая из которых задаёт 
конкретную задачу управления, позволяющую оценить способность метода 
корректировать ошибки спецификации награды и улучшать поведение агента. 
Во всех случаях истинная функция вознаграждения недоступна в реальных задачах, 
но используется авторами для генерации эталонных предпочтений.

\begin{itemize}
    \item \emph{Среда управления пандемией (Pandemic Mitigation, модель SEIR)} 
    \cite{kompella2020covid}~-- агент должен выбирать интенсивность 
    ограничительных мер, стремясь уменьшить распространение инфекции 
    при минимальных социально-экономических потерях.

    \item \emph{Среда контроля уровня глюкозы (Glucose Monitoring, симулятор UVA/Padova)} 
    \cite{man2014uvapadova}~-- агент регулирует дозировку инсулина, удерживая 
    уровень глюкозы пациента в безопасном диапазоне.

    \item \emph{Среда управления дорожным движением (Traffic Control, FLOW)}
    \cite{wu2021flow}~-- агент управляет транспортным средством или частью 
    транспорта, стремясь снижать заторы и увеличивать пропускную способность 
    дорожной сети.

    \item \emph{Среда задач безопасного поведения (AI Safety Gridworlds)} 
    \cite{leike2017aisg}~-- агент должен избегать ловушек награды и 
    нежелательных стратегий в специально сконструированных дискретных мирах, 
    моделирующих ошибки спецификации целей.
\end{itemize}

Траектории для генерации предпочтений формируются с использованием истинной 
функции вознаграждения. Тем самым предпочтения служат единственным источником 
информации об ошибках прокси-награды, позволяя объективно оценить работу метода.

\subsection{Метрики}

Основная метрика представляет собой \emph{среднее истинное вознаграждение}, 
получаемое агентом при взаимодействии со средой. Эта величина показывает, 
насколько поведение агента соответствует фактическим целям и намерениям, 
заложенным в истинной награде.

Дополнительно анализируются две характеристики динамики обучения.
\begin{itemize}
    \item  \emph{устойчивость обучения}, оцениваемая по изменению качества 
политики при увеличении числа предпочтений: метод должен демонстрировать 
монотонное или стабильное улучшение без резких деградаций.
    \item \emph{восстановление после искажения вознаграждения}, то есть 
    способность агента вернуться к корректному поведению, если исходная 
    прокси-награда стимулирует нежелательные действия.
\end{itemize} 


\subsection{Модели}

Политика агента обучается с помощью метода \emph{оптимизации проксимальных 
политик, PPO} \cite{schulman2017ppo}, обеспечивающего стабильное обновление 
параметров при оптимизации модифицированной функции награды. Этот же алгоритм 
используется во всех средах для получения поведения, на основе которого 
формируются траектории для сравнения.

Ссылочные политики $\pi_{\text{ref}}$ получаются методом 
\emph{имитации поведения}, при котором модель обучается воспроизводить 
заранее заданные демонстрации. Для построения таких политик используются 
подходы, описанные в \cite{treiber2000idm} и \cite{willis1999pid}.

